% JAC1964-20261121, /root/arithmetic_sign.
\subsection*{AR-28. The exact nef range of the canonical diagonal metric}

\begin{proposition}[AR-28]
Let $C/K$ be a smooth projective geometrically integral curve of genus
$g\ge2$ over a number field.  Suppose a degree-one line bundle
$\alpha$ with $(2g-2)\alpha=\omega_C$ is defined over $K$;
in particular one may take $\alpha=\mathcal O(e)$ for the canonical
point supplied by AR-27.  Give it the metric
$\overline\alpha=(2g-2)^{-1}\overline\omega_a$, and give the
diagonal its canonical admissible metric $\overline\Delta_a$.
For $\delta_1,\delta_2,t\in\mathbf Q$, put
\[
 \overline D_{\delta_1,\delta_2,t}
 =\delta_1\overline\alpha_1+\delta_2\overline\alpha_2
       +t(\overline\Delta_a-\overline\alpha_1
                              -\overline\alpha_2).
\]
Then, for these specified adelic metrics,
\[
 \overline D_{\delta_1,\delta_2,t}\text{ is nef}
 \quad\Longleftrightarrow\quad
 \delta_1\ge0,\quad\delta_2\ge0,\quad
 \delta_1\delta_2\ge g^2t^2.                            \tag{AR.28a}
\]
This is an adelic nefness criterion; it is not asserted to be the
geometric nef cone of the underlying divisor classes.

More precisely, let $J=\operatorname{Jac}(C)$, let
$i_\alpha(x)=[x-\alpha]$, and let $\overline\theta$ be the
canonically metrized rational principal polarization on $J$.
Write
\[
 A_j=p_j^*i_\alpha^*\overline\theta,
 \qquad
 \overline c=\langle\overline\alpha,\overline\alpha\rangle_{C/K},
\]
and let $Q$ be the pullback of
$m^*\overline\theta-p_1^*\overline\theta-p_2^*\overline\theta$
under $i_\alpha\times i_\alpha$.  As rational adelic line bundles,
\[
 \overline D_{\delta_1,\delta_2,t}
   =\frac{\delta_1}{g}A_1+\frac{\delta_2}{g}A_2-tQ
      +\left(\frac{\delta_1+\delta_2}{2g}-t\right)\pi^*\overline c.
                                                               \tag{AR.28b}
\]
When (AR.28a) holds, the first three terms are a canonical positive
semidefinite pullback from $J^2$, and the last term is a nef base
line with a nonnegative coefficient.
\end{proposition}

\begin{proof}
We use Yuan's exact metric identities with $\Theta=2\theta$.
His Theorem 2.9(1), for a degree-one admissible $\alpha$, and
Theorem 2.10(2), for $j(x,y)=y-x$, give
\[
 i_\alpha^*\overline\theta
       =g\overline\alpha-\tfrac12\pi^*\overline c,
 \qquad
 j^*\overline\theta
       =\overline\Delta_a+(g-1)
                    (\overline\alpha_1+\overline\alpha_2).
                                                               \tag{AR.28c}
\]
These equalities include the metrics, not just their curvatures or
underlying line bundles.  Since
$j^*\overline\theta=A_1+A_2-Q$, they imply
\[
 Q=\overline\alpha_1+\overline\alpha_2
                  -\overline\Delta_a-\pi^*\overline c.
\]
Substitution proves (AR.28b), including its base-line constant.
The line $\overline c$ is nef: $\overline\omega_a$ is nef by
Yuan's Theorem 2.10(1), and the Deligne self-pairing of the nef
line $\overline\alpha$ is nef on $\operatorname{Spec}K$.

If the numerical conditions in (AR.28a) hold, the matrix
\[
 H=\begin{pmatrix}\delta_1/g&-t\\-t&\delta_2/g\end{pmatrix}
\]
is positive semidefinite.  It defines a canonical nef rational
line bundle on $J^2$.  One may verify this directly by expressing
the rational positive semidefinite binary form as a sum of
nonnegative rational multiples of squares of rational linear
forms; the corresponding bundles are pullbacks of
$\overline\theta$ under rational homomorphisms $J^2\to J$,
with denominators cleared.  Its pullback to $C^2$ is the first
three terms of (AR.28b).  Also
\[
 \frac{\delta_1+\delta_2}{2g}-t
       \ge\frac{\sqrt{\delta_1\delta_2}}{g}-t
       \ge |t|-t\ge0.
\]
The last term is therefore nef as well.  This proves sufficiency
at every place and for all arithmetic point heights, rather than
only complex semipositivity.

Conversely, nefness restricts to the two geometric fibers of
$C^2$, whose degrees are $\delta_1$ and $\delta_2$; thus both
are nonnegative.  At any complex embedding let $\mu$ be the
Arakelov probability form, so that
$c_1(\overline\alpha)=\mu$ and
$c_1(i_\alpha^*\overline\theta)=g\mu$ by (AR.28c).
At a diagonal point $(x,x)$, use the same local parameter on the
two curve factors.  The differential of $j$ is the difference of
the two differentials of $i_\alpha$.  Therefore the Hermitian
matrix of $c_1(j^*\overline\theta)$ there is
\[
 g\mu(x)\begin{pmatrix}1&-1\\-1&1\end{pmatrix}.
\]
The second identity in (AR.28c) then gives for
$c_1(\overline\Delta_a)$ the matrix
\[
 \mu(x)\begin{pmatrix}1&-g\\-g&1\end{pmatrix}.
\]
Consequently the curvature matrix of the stated metric on
$\overline D_{\delta_1,\delta_2,t}$ at $(x,x)$ is
\[
 \mu(x)\begin{pmatrix}\delta_1&-gt\\-gt&\delta_2\end{pmatrix}.
                                                               \tag{AR.28d}
\]
The form $\mu(x)$ is strictly positive, since the canonical
linear system has no base point.  An adelic nef metric is
semipositive at this complex place.  Its determinant must
therefore be nonnegative, proving
$\delta_1\delta_2\ge g^2t^2$.  This proves necessity.
\end{proof}

\paragraph{The point metric when $\alpha=\mathcal O(e)$.}
The equality of underlying line bundles does not identify
$\overline\alpha$ with the canonical point metric
$\overline{\mathcal O(e)}_a$.  In fact the exact identities are
\[
 e^*\overline\alpha=\frac1{2g}\overline c,
 \qquad
 \overline{\mathcal O(e)}_a
      =\overline\alpha-\frac{2g-1}{2g}\pi^*\overline c.
                                                               \tag{AR.28f}
\]
The first follows by evaluating the first equality of (AR.28c)
at $e$, since $i_\alpha(e)=0$ and the polarization is rigidified.
For the second, both metrics are admissible on the same degree-one
line and hence differ by a base line.  Admissible adjunction gives
$e^*\overline{\mathcal O(e)}_a=-e^*\overline\omega_a$,
which determines that base line as displayed.  Thus the normalized
degree-zero diagonal correspondence with its axes rigidified at
$e$ is
\[
 \overline t_e
 =\overline\Delta_a-p_1^*\overline{\mathcal O(e)}_a
                    -p_2^*\overline{\mathcal O(e)}_a
                    +\pi^*e^*\overline{\mathcal O(e)}_a
 =\overline\Delta_a-\overline\alpha_1-\overline\alpha_2
                    +\pi^*\overline c=-Q.
\]
These constants keep the canonical-point normalization of AR-27
compatible with the metrized tensor correspondences in AR-25.

\paragraph{What the Vojta small-section theorem does and does not supply.}
Geometrically, the same divisor has
\[
 D_{\delta_1,\delta_2,t}^{2}
       =2(\delta_1\delta_2-gt^2),                       \tag{AR.28e}
\]
because $(\Delta-\alpha_1-\alpha_2)^2=-2g$ and this last
class is orthogonal to each $\alpha_j$.  With
$\delta_1,\delta_2>0$, the interval
\[
 gt^2<\delta_1\delta_2<g^2t^2
\]
therefore gives big underlying divisors but non-nef canonical
metrics.  Bigness follows from surface Riemann--Roch and positive
intersection with $\alpha_1+\alpha_2$; it is not a nefness claim.
The cited Vojta argument chooses the product just above $gt^2$,
inside this range.
Adding any adelic base line leaves the curvature (AR.28d)
unchanged, so it cannot make these metrics nef.

This precisely matches the primary proof of Yu--Yuan--Zhou,
\emph{Quantitativity on the number of rational points in the
Mordell conjecture}, Proposition 3.7, author pp.~37--38.  There
the Vojta bundle is a scalar multiple of this $D$; a positive
base shift is used to make its arithmetic volume positive, by
writing it as a difference of two nef bundles and applying the
arithmetic Siu inequality.  The conclusion is a small section
with an explicit error term.  It does not say that the shifted
Vojta bundle is nef, and so it cannot be inserted as the nef
polarization in Yuan--Zhang's Hodge-index inequality.

\paragraph{Scope for the remaining complementary height.}
Formula (AR.28b) represents every nef member of this specified
diagonal/canonical family by a Jacobian positive semidefinite
bundle and a nonnegative base term.  The Jacobian $J(C)$ may be
larger than the fixed $E^3$ quotient.  Thus this representation
does not assert that all resulting mixed identities were already
covered by AR-26; its complementary Jacobian factors must still
be tracked.  Nor does it replace the projected modified diagonal
of AR-27 by its total height.  The result eliminates the proposed
use of the small-section theorem as an unshifted nefness theorem,
while fixing all metric constants for further diagonal tests.

\paragraph{Primary locators.}
Yuan, \emph{Arithmetic bigness and a uniform Bogomolov-type
result}, author version dated 21 August 2024,
Theorem 2.9(1), pp.~35--36, and Theorem 2.10(1)--(2),
pp.~37--38, read directly; published in Annals of Mathematics
203 (2026), pp.~15--119.  The author's PDF uses $\Theta$ for
twice the principal polarization, which is why it is divided
by two throughout (AR.28c).  The normalization of admissible
metrics is also recorded in Yu--Yuan--Zhou, \S2.5, pp.~29--33.
